\providecommand{\itescastudentname}{Martin Jonathan de la Cruz}
\providecommand{\itescastudentid}{26130503}
\providecommand{\itescastudentemail}{26130503@cajeme.tecnm.mx}
\providecommand{\itescafaculty}{Maestría en Gestión Administrativa}
\providecommand{\itescacoursename}{Seminario I}
\providecommand{\itescacoursecode}{MGA-SEMINARIO-I}
\providecommand{\itescadepartment}{Posgrado e Investigación}
\providecommand{\itescateachingfigure}{Dra. Carla Olimpya Zapuche Moreno}
\providecommand{\itescaprogramlevel}{Maestría}
\providecommand{\itescaprogramperiod}{Agosto--noviembre de 2026}
\providecommand{\itescamodality}{En línea}
\providecommand{\itescalocation}{Cajeme, Sonora}
\providecommand{\itescaasset}{ITESCA/assets-itesca/itesca-monograma.png}
\providecommand{\itescaofficialasset}{ITESCA/assets-itesca/web/logo-itesca-contorno.png}
\providecommand{\itescasecondaryasset}{ITESCA/assets-itesca/itesca-monograma.png}
\providecommand{\itescabibliography}{ITESCA/maestria-en-gestion-administrativa/seminario-i-mga/seminario-i}

\def\itescadocumenttitle {Mi primera práctica en formato APA}
\def\itescadocumentsubtitle {Actividad 2}
\def\itescadocumentsubject {Aplicación de citas y referencias APA 7}
\def\itescaactivityname {Actividad 2. Mi primera práctica en formato APA}
\def\itescasessionname {Unidad 2}
\def\itescablockname {Seminario I}
\def\itescamainproduct {Práctica de citación y referencias APA}

\documentclass[spanish,letterpaper,oneside]{article}

\def\documenttitle {\itescadocumenttitle}
\def\documentsubtitle {\itescadocumentsubtitle}
\def\documentsubject {\itescadocumentsubject}
\def\documentauthor {\itescastudentname}
\def\coursename {\itescacoursename}
\def\coursecode {\itescacoursecode}
\def\universityname {Instituto Tecnológico Superior de Cajeme}
\def\universityfaculty {\itescafaculty}
\def\universitydepartment {\itescadepartment}
\def\universitydepartmentimage {\itescaofficialasset}
\def\universitydepartmentimagecfg {width=5.2cm}
\def\universitylocation {\itescalocation}

\def\authortable {
  \begin{tabular}{ll}
    \textbf{Alumno:} & \itescastudentname \\
    No. de control: & \itescastudentid \\
    Carrera: & \itescafaculty \\
    Semestre: & \itescaprogramperiod \\
    Asignatura: & \itescacoursename \\
    Docente: & \itescateachingfigure \\
    Modalidad: & \itescamodality \\
    Fecha de entrega: & 30 de agosto de 2026 \\
    Lugar: & \universitylocation
  \end{tabular}
}

\input{ITESCA/maestria-en-gestion-administrativa/seminario-i-mga/assets-seminario-i/plantilla/template.tex}
\usepackage{helvet}
\renewcommand{\familydefault}{\sfdefault}
\usepackage{ragged2e}
\usepackage{enumitem}
\usepackage{tikz}
\usetikzlibrary{calc}

\setlength{\parindent}{1.27cm}
\setlength{\parskip}{0pt}

\newenvironment{apablockquote}{%
  \begin{list}{}{\setlength{\leftmargin}{1.27cm}\setlength{\rightmargin}{0cm}}\item\relax
}{\end{list}}

\newenvironment{apareferences}{%
  \begin{list}{}{\setlength{\leftmargin}{1.27cm}%
    \setlength{\itemindent}{-1.27cm}%
    \setlength{\itemsep}{0.8\baselineskip}%
    \setlength{\parsep}{0pt}}}{\end{list}}

\definecolor{itescaRedDark}{HTML}{8F1117}
\definecolor{itescaRed}{HTML}{D70F18}
\definecolor{itescaCharcoal}{HTML}{141414}
\definecolor{itescaSilver}{HTML}{D9D9D9}
\def\portraittitlecolor {itescaRedDark}

\newcommand{\insertcoverwatermark}{%
  \AddToShipoutPictureBG*{%
    \begin{tikzpicture}[remember picture,overlay]
      \node[opacity=0.13,inner sep=0pt,yshift=-0.10cm] at (current page.center) {\includegraphics[width=0.58\paperwidth]{\itescaasset}};
    \end{tikzpicture}%
  }%
}
\newcommand{\insertcoverinstitutionalfooter}{%
  \AddToShipoutPictureFG*{%
    \begin{tikzpicture}[remember picture,overlay]
      \fill[itescaCharcoal] (current page.south west) rectangle ([yshift=1.35cm]current page.south east);
      \fill[itescaRed] ([yshift=1.35cm]current page.south west) rectangle ([yshift=1.47cm]current page.south east);
      \node[anchor=south west] at ([xshift=0.95cm,yshift=0.16cm]current page.south west) {\includegraphics[height=0.88cm]{\itescasecondaryasset}};
      \node[anchor=south east,text=white,align=right,font=\sffamily\bfseries\small] at ([xshift=-0.95cm,yshift=0.28cm]current page.south east) {Instituto Tecnológico Superior de Cajeme\\\textcolor{itescaSilver}{\itescalocation\ \textperiodcentered\ itesca.edu.mx}};
    \end{tikzpicture}%
  }%
}

\begin{document}
\insertcoverwatermark
\insertcoverinstitutionalfooter
\templatePortrait
\templatePagecfg

\begin{abstractd}
\textbf{Objetivo:} Aplicar las normas APA (7.\textsuperscript{a} edición) en la presentación formal, las citas dentro del texto y la lista de referencias de un escrito breve relacionado con la gestión administrativa.

\textbf{Producto:} Práctica académica sobre gestión administrativa basada en evidencia que integra una paráfrasis, una cita textual corta, una cita textual larga y seis fuentes de cinco tipos documentales.
\end{abstractd}

\templateIndex
\templateFinalcfg
\doublespacing
\setlength{\parindent}{1.27cm}

\section{Introducción}

La toma de decisiones constituye un problema central de la gestión administrativa porque una alternativa puede parecer razonable y, aun así, carecer de datos verificables o atribuir como propia una idea ajena. El uso de APA 7 permite diferenciar con precisión la voz del autor y la evidencia consultada. El objetivo es aplicar las reglas de citación y referencia en un argumento coherente sobre decisiones administrativas basadas en evidencia, mediante una paráfrasis, una cita breve y una cita en bloque vinculadas con fuentes recuperables.

\section{Evidencia y decisiones en la gestión administrativa}

\subsection{Paráfrasis y cita textual corta}

La gestión administrativa basada en evidencia permite sustituir decisiones apoyadas únicamente en la intuición por juicios que integran datos de la organización, conocimiento científico y experiencia profesional. Rousseau (2006) sostiene que este enfoque convierte los mejores hechos disponibles en una base para las prácticas directivas; en el mismo sentido, Hernández Sampieri et al. (2014) explican que la investigación sistemática articula procesos críticos y empíricos para estudiar problemas, por lo que sus resultados pueden orientar diagnósticos y decisiones sin confundir una opinión con una conclusión sustentada. Esta perspectiva es especialmente útil en mercadotecnia y ventas: antes de modificar una campaña, un proceso o una meta, la persona responsable debe formular una pregunta, comparar indicadores, valorar la calidad de las fuentes y documentar el criterio elegido. La complejidad cotidiana no desaparece por seguir una secuencia formal; como advierten Bright y Cortes (2019), el mundo real de quienes administran es un \textquote{flujo desordenado y frenético de actividad continua} (cap. 1, párr. 4, traducción propia). La cita corta resume una dificultad concreta, mientras que la paráfrasis anterior integra ideas ajenas con una redacción propia y conserva la atribución correspondiente.

\subsection{Cita textual larga}

La necesidad de interpretar la evidencia dentro de su contexto también se aprecia en la siguiente descripción del trabajo directivo:

\begin{apablockquote}
La mayoría de los libros de administración dirían, como lo hace este, que los gerentes dedican su tiempo a planear, organizar, dotar de personal, dirigir, coordinar, informar y controlar. Estas actividades, como encontró Hannaway en su estudio de los gerentes en el trabajo, \textquote{no describen, de hecho, lo que hacen los gerentes}. En el mejor de los casos parecen describir objetivos vagos que los gerentes intentan alcanzar continuamente. (Bright \& Cortes, 2019, cap. 1, párr. 4, traducción propia)
\end{apablockquote}

El pasaje supera 40 palabras y, por ello, se presenta como cita en bloque, sin comillas externas, con sangría y localizador. Su contenido muestra que una lista de funciones no basta para explicar una decisión administrativa: es necesario reconstruir el problema, las restricciones y los efectos esperados. Un antecedente académico aplicado examinó la toma de decisiones en la gestión administrativa de una empresa (Zambrano Andrade, 2018), mientras que Sinek (2009) propuso comunicar primero el propósito para movilizar la acción. Ambos recursos cumplen funciones distintas: una tesis aporta un estudio documentado y el video ofrece una exposición profesional que debe evaluarse de acuerdo con su naturaleza. La American Psychological Association (2022) indica que las citas menores de 40 palabras se incorporan entre comillas y que las de 40 palabras o más se disponen en bloque.

\section{Conclusiones}

Considero que el uso correcto de APA no es un adorno tipográfico, sino un mecanismo de integridad académica. Mi postura se sustenta en que una decisión administrativa solo puede someterse a crítica cuando el lector distingue qué parte corresponde al análisis propio y cuál procede de una fuente verificable. En consecuencia, citar con precisión fortalece la argumentación, reduce atribuciones ambiguas y establece una práctica transferible a la futura tesis y a los informes profesionales.\footnote{Se utilizó inteligencia artificial como apoyo para organizar ideas y revisar la redacción; la selección de fuentes, su verificación y el análisis final fueron realizados y supervisados por el autor.}

\clearpage
\section*{Referencias}
\addcontentsline{toc}{section}{Referencias}
\begin{apareferences}
  \item American Psychological Association. (2022, julio). \emph{Quotations}. APA Style. \url{https://apastyle.apa.org/style-grammar-guidelines/citations/quotations}

  \item Bright, D. S., \& Cortes, A. H. (2019). \emph{Principles of management}. OpenStax, Rice University. \url{https://openstax.org/details/books/principles-management}

  \item Hernández Sampieri, R., Fernández Collado, C., \& Baptista Lucio, P. (2014). \emph{Metodología de la investigación} (6.\textsuperscript{a} ed.). McGraw-Hill Education.

  \item Rousseau, D. M. (2006). Is there such a thing as ``evidence-based management''? \emph{Academy of Management Review, 31}(2), 256--269. \url{https://doi.org/10.5465/amr.2006.20208679}

  \item Sinek, S. (2009, septiembre). \emph{How great leaders inspire action} [Video]. TED Conferences. \url{https://www.ted.com/talks/simon_sinek_how_great_leaders_inspire_action}

  \item Zambrano Andrade, L. F. (2018). \emph{La toma de decisiones en la gestión administrativa de la empresa Arsaico Cía. Ltda., cantón Chambo, período 2016} [Tesis de licenciatura, Universidad Nacional de Chimborazo]. Repositorio Digital UNACH. \url{http://dspace.unach.edu.ec/handle/51000/4584}
\end{apareferences}

\end{document}
